\section{Limitations and Conclusion}

\textbf{Data limitations.} Fashion-MNIST contains centered, low-resolution grayscale objects with balanced labels. Results may not transfer to color images, cluttered backgrounds, larger resolutions, or naturally imbalanced data. Several upper-body categories are intrinsically ambiguous at $28\times28$ resolution.

\textbf{Method limitations.} All five architectures were implemented and evaluated, but the configurations were fixed rather than exhaustively tuned and model capacities are not matched. No formal augmentation, dropout, representation, or parameter-budget ablation was run. ViT's four-layer encoder and 50-token sequence are computationally expensive for this small input, while the model's learned global attention does not encode locality as strongly as a CNN.

\textbf{Experimental limitations.} Results come from one deterministic seed, so run-to-run variance is unknown. The exact CPU model and complete session identity were not recorded; therefore, training and inference measurements indicate practical cost but are not hardware-controlled benchmarks. The qualitative set is deliberately small and illustrates behavior rather than estimating frequency.

\textbf{Expected generalization and future work.} CNN's advantage is observed on this dataset and configuration; its magnitude and ranking should be confirmed across seeds and datasets. Future work is to repeat experiments across multiple seeds, log exact hardware and software automatically.

\textbf{Conclusion.} Under the shared Fashion-MNIST pipeline, architectural inductive bias materially changes accuracy and efficiency across all five models. Linear is the cheapest but underfits; MLP gains from nonlinear capacity; GRU provides a strong accuracy--parameter compromise; and CNN is the strongest measured classifier, reaching 92.55\% accuracy and 0.9259 macro-F1. ViT reaches 87.93\% accuracy and 0.8771 macro-F1 but has the highest measured inference time and FLOPs. Thus, global patch attention does not outperform the convolutional or row-recurrent alternatives in this experiment. CNN is preferred when predictive quality is primary, Linear when minimal size and latency matter, and GRU when parameter efficiency matters. These conclusions are limited to one seed and incompletely characterized CPU hardware.
